% TODO: 目前内容不足以凑满两页（正反 A4 纸），所以有些内容注释掉了，之后完善

\documentclass{article}
\setlength\parindent{0pt}
\pagenumbering{gobble}

\usepackage{geometry}
\geometry{papersize={210mm,297mm}}
\geometry{left=9mm,right=9mm,top=9mm,bottom=9mm}

\usepackage{fontspec}
\setmainfont{DejaVu Sans}[Scale=0.9]

\usepackage{xeCJK}
\CJKfamily{zhsong}

\usepackage[svgcolors]{xcolor}
\usepackage{longfbox}
\usepackage{epstopdf}
\usepackage{float}
\usepackage{xpatch}
\usepackage{tipa}
\usepackage[export]{adjustbox}
\usepackage{enumitem}
\usepackage{mdframed}

\usepackage{multirow}
\usepackage{tabularx}
\renewcommand\tabularxcolumn[1]{m{#1}} % for vertical centering text in X column

\usepackage{multicol}
\usepackage[most]{tcolorbox}
\usepackage{amssymb}
\setlength{\columnsep}{3mm}

\usepackage[explicit,compact]{titlesec}
\titleformat{\section}{\normalfont\bfseries}{}{0pt}{【#1】}

% 啊！万能的 StackExchange！
% https://tex.stackexchange.com/questions/475466/latex-three-column-layout-merging-two-of-them-at-the-begining
%
\newlength{\abstractwidth}
\newlength{\columnshrink}
\newsavebox{\twocolinsert}
%
\makeatletter
\newlength{\resized@col}
\newcounter{column@count}
%
\xpatchcmd{\multi@column@out}{
	\process@cols\mult@gfirstbox{%
		\setbox\count@
		\vsplit\@cclv to\dimen@
		\set@keptmarks
		\setbox\count@
		\vbox to\dimen@
		{\unvbox\count@ \ifshr@nking\vfilmaxdepth\fi}%
	}%
}{
	\process@cols\mult@gfirstbox{%
		\global\advance\c@column@count\@ne
		\resized@col\dimen@%
		\ifnum\c@column@count=\tw@
				\advance\resized@col-\columnshrink
		\fi%
		\setbox\count@
		\vsplit\@cclv to\resized@col
		\set@keptmarks
		\setbox\count@
		\vbox to\dimen@{
			\ifnum
				\c@column@count=\tw@ \vspace*{\columnshrink}
			\fi
			\unvbox\count@
			\ifshr@nking\vfilmaxdepth\fi
		}%
	}%
}{\typeout{Success}}{\typeout{Failure}}
\makeatother

% for designing header
\newsavebox\mysavebox
\newenvironment{imgminipage}[2][]{%
   \def\imgcmd{\includegraphics[width=\wd\mysavebox, height=\dimexpr\ht\mysavebox+\dp\mysavebox\relax, #1]{#2}}%
   \begin{lrbox}{\mysavebox}%
   \begin{minipage}%
}{%
   \end{minipage}
   \end{lrbox}%
   \sbox\mysavebox{\setlength{\fboxrule}{0pt}\fbox{\usebox\mysavebox}}%
   \mbox{\rlap{\raisebox{-\dp\mysavebox}{\imgcmd}}\usebox\mysavebox}%
}

\renewcommand{\labelitemi}{$\blacktriangleright$}

\tcbset{
    frame code={}
    center title,
    left=0pt,
    right=0pt,
    top=6pt,
    bottom=0pt,
    colback=gray!40,
    colframe=white,
    enlarge left by=0mm,
    boxsep=0pt,
    arc=0pt,outer arc=0pt,
}

\begin{document}
\begin{multicols*}{3}

	\setlength{\abstractwidth}{2\linewidth}
	\addtolength{\abstractwidth}{\columnsep}
	\savebox{\twocolinsert}{
	\begin{minipage}{\abstractwidth}
		\noindent 核准日期：2026年06月19日
		\newline 编辑日期：2026年06月18日
		\newline

		 \begin{mdframed}[leftline=false, rightline=false, innertopmargin=0pt, innerbottommargin=0pt, innerrightmargin=0pt, innerleftmargin=2em]
		 	\includegraphics[width=0.15\abstractwidth, valign=m]{assets/debian-text.eps}
		 	\hfill
		 	\begin{imgminipage}{assets/header-background.eps}[t]{0.7\abstractwidth}
		 		\Large \textbf{盒装安装媒介说明书}
		 		\normalsize 请仔细阅读说明书并在管理员指导下使用
		 	\end{imgminipage}
		 \end{mdframed}

		\includegraphics[width=\abstractwidth]{assets/header.eps}


		\begin{mdframed}[hidealllines=true, innerbottommargin=.5em, innertopmargin=0pt]
			\sffamily

			{\centering 警示语 \par}

			无论是否与其它操作系统合用，安装 Nyarch linux 均存在丢失磁盘上所有内容的风险。（参见【不良反应】）

            本说明中的信息和意见无意构成法律建议，法律建议可通过咨询律师来获得。

			一些硬件制造商拒绝告诉我们如何给他们的硬件编写驱动程序。另一些则要求签署不公开的契约才能接触相关文档，以阻止我们以自由软件形式发布驱动程序源代码。由于我们未被授权使用这些文档，因此它们无法在 Nyarch linux 下工作。
		\end{mdframed}
	\end{minipage}}
	\setlength{\columnshrink}{\ht\twocolinsert}
	\addtolength{\columnshrink}{\dp\twocolinsert}
	\noindent\usebox{\twocolinsert}


	\begin{tcolorbox}
	\section*{发行版名称}
	\end{tcolorbox}
	\begin{tabularx}{\linewidth}{@{}ll@{}}
		通用名称： & Nyarch linux \\
		正式名称： & Nyarch linux \\
		英文音标： & \textipa{/ˈnjaːtʃ ˈlɪn.əks/}
	\end{tabularx}

	\medskip


	\begin{tcolorbox}
	\section*{内容}
	\end{tcolorbox}

	Nyarch Linux 是一个基于Arch Linux的为堕落的御宅族提供的完美发行版。通过滚动更新尽可能地提供最新的稳定版软件的同时，提供了每个御宅爱好者都必须预装的最佳开源应用。

	内核版本：26.04

	\medskip


	\begin{tcolorbox}
	\section*{性质}
	\end{tcolorbox}

	本系统为采用 GNU/Linux 内核的操作系统，安装后可由 UEFI 或 Legacy 引导方式启动。

	\medskip


	\begin{tcolorbox}
	\section*{适应平台}
	\end{tcolorbox}

	\begin{itemize}
		\item 仅支持使用 x86-64 (AMD64) 架构的计算机、服务器和嵌入式设备。
		\item 本包装盒中的安装媒介适用于何平台以实际为准。
	\end{itemize}


	\begin{tcolorbox}
	\section*{规格}
	\end{tcolorbox}

	1 枚刻录有 Nyarch Linux 26.04 安装映像的可引导 USB 闪存驱动器。

	\medskip

	\begin{tcolorbox}
	\section*{用法}
	\end{tcolorbox}

	使用 USB 设备引导并启动你的设备。

	访问设备的 BIOS/UEFI 工具。访问方式取决于你的设备。 如果你不知道如何访问 BIOS，在设备启动时狂按 F 键并大喊Hora Hora Hora，直到按对为止。
	找到启动顺序设置，将你的 U 盘放在第一位。或者，如果你有访问启动菜单的权限，从那里启动。最后，保存并退出。

	随后，你会进入 Live 环境中，并且自动打开 Nyarch 安装程序，然后它会提示你输入本地化设置，选择如何安装。根据你的偏好填写，然后点击确定，等待安装完成，然后重启！
	\medskip	

	\begin{tcolorbox}
	\section*{不良反应}
	\end{tcolorbox}

	Nyarch linux 有一个对用户和开发者所提交的软件缺陷报告进行归档管理的缺陷跟踪系统，每个软件缺陷报告都被授予一个编号并且被长期跟踪，直到它被标记为已修复。
	此外， Nyarch 团队规模较小，除了正常的Arch滚动发布风险外，你可能会遇到项目特定的bug。

	\medskip


	\begin{tcolorbox}
	\section*{注意事项}
	\end{tcolorbox}
	\begin{itemize}[leftmargin=*]

		\item 满足最低的硬件要求

		下表是内存和硬盘的需求。

		{\small\begin{tabularx}{\linewidth}{|X|X|X|X|}
			\hline
			类别 & RAM\newline (最低) & RAM\newline (推荐) & 硬盘 \\
			\hline
			Gnome & 3GB & 6GB & 30GB \\
			\hline
			KDE plasma & 3GB & 6GB & 30GB \\
			\hline
		\end{tabularx}}

		基于您的需求，也许可以使用低于上表所列的配置完成系统安装。但是多数用户在无视这些建议的情况下会安装失败。

		\item 需要固件的设备

		除了需要设备驱动程序，有些硬件还要在使用之前加载固件（firmware）或微码（microcode）。

	\end{itemize}


	\begin{tcolorbox}
	\section*{禁忌}
	\end{tcolorbox}

	在操作过程中出现或即将出现下列任何一种情况，请立即停止操作，并准备好系统恢复。

	\begin{itemize}[leftmargin=*]
		\setlength{\itemsep}{0pt}
		\setlength{\parskip}{0pt}
		\setlength{\parsep}{0pt}

		\item 以根权限在根目录下执行递归删除
		\item 未确认设备名即使用 dd 命令
		\item 未确认设备名即执行格式化
		\item 未确认操作即使用命令重定向
		\item 未经确认即执行来自网络的脚本
		\item 长期在散热不良的设备上高负载使用
	\end{itemize}

    % 注释部分就不管了

	 \begin{tcolorbox}
	 \section*{版本迭代}
	 \end{tcolorbox}
	 	\begin{itemize}[leftmargin=*]
		\setlength{\itemsep}{0pt}
		\setlength{\parskip}{0pt}
		\setlength{\parsep}{0pt}

		\item Nyarch 会不定期进行版本迭代更新，仅需前往 Nyarch Updater 更新即可。

	\end{itemize}


	\begin{tcolorbox}
	\section*{系统相互作用}
	\end{tcolorbox}
	\begin{itemize}[leftmargin=*]
		\setlength{\parindent}{0pt}

		\item 与 Windows 的相互作用

		当您有双引导时，若另一个操作系统与 Windows 访问相同的文件系统，这就有可能会导致问题和数据丢失。在这种情况下，文件系统的真实状态可能与 Windows 认为在“启动”之后的情况不同，并且可能在进一步写入文件系统时导致文件系统损坏。因此，在双引导设置中，为了避免文件系统损坏，有必要在 Windows 中禁用“快速启动”功能。

		在罕见情况中已观察到，在使用 Windows 进行系统更新时，可能会出现重新启动后引导被破坏从而导致 Nyarch linux 无法启动的情况。同时，若在安装过程中将引导信息写入与 Windows 所在的物理磁盘 MBR 内，将导致后者无法正常启动。

		\item 与其他 Linux 发行版的相互作用

		与 Manjaro 混用时，可能导致 AUR 宕机。

	\end{itemize}


	\begin{tcolorbox}
	\section*{贮藏}
	\end{tcolorbox}

	-40℃\textasciitilde +70℃

	妥善贮藏所有安装媒介，勿使非安装人员触及。

	\medskip


	\begin{tcolorbox}
	\section*{包装}
	\end{tcolorbox}

	刻录有 Nyarch linux 26.04 版本安装映像的、兼容 USB 3.0/2.0 协议的闪存驱动器。

	1 枚/盒

	\medskip


	\begin{tcolorbox}
	\section*{有效期}
	\end{tcolorbox}

	Nyarch linux 使用半滚动升级模式，普通软件更新只需执行一条命令，便可升级软件到最新版本。而 Nyarch 的版本更新则需要依靠 Nyarch Updater 。 基于上述原因，对 Nyarch linux 的支持长期有效。

	\medskip


	\begin{tcolorbox}
	\section*{执行标准}
	\end{tcolorbox}
	\begin{tabularx}{\linewidth}{@{}ll@{}}
		\multirow{4}{*}{}{开源许可证：} & GNU GPL（主要）\\
		~ & GNU LGPL \\
		~ & BSD \\
		~ & 以及其他授权条款 \\
	\end{tabularx}

	\medskip


	\begin{tcolorbox}
	\section*{批准文号}
	\end{tcolorbox}

	说明书使用 CC-BY-SA 3.0 协议授权。

	\medskip


% 	\begin{tcolorbox}
% 	\section*{生产单位}
% 	\end{tcolorbox}
%
% 	Debian 计划
%
% 	\medskip


	\begin{tcolorbox}
	\section*{说明书}
	\end{tcolorbox}
	\begin{tabularx}{\linewidth}{@{}ll@{}}
		\multirow{2}{*}{}{编审：} & @Skywind_Fox\\
		翻译：@Skywind_Fox\\
		翻译校对：
		排版: & @Skywind_Fox\\
		GitHub: & Skywindfox/Nyarch-media-box\\
	\end{tabularx}

	\medskip


	\vfill
	\begin{flushright}
		Nyarch linux\linebreak 20260618
		\linebreak
		\newline
	\end{flushright}

\end{multicols*}
\end{document}